\documentclass[conference]{IEEEtran}
\usepackage{amsmath,amssymb,amsfonts}
\usepackage{algorithmic}
\usepackage{graphicx}
\usepackage{textcomp}
\usepackage{xcolor}
\usepackage{booktabs}
\usepackage{cite}
\usepackage{microtype}
\usepackage{tabularx}
\usepackage{url}
\usepackage{balance}

\begin{document}

\title{Architectural Deep Dive into Qwen2.5-72B: Design Foundations, Attention Dynamics, and Scalable Inference in Dense Frontier Language Models}

\maketitle

\begin{abstract}
Dense autoregressive Transformers remain the foundational backbone of modern artificial intelligence. Among open-weights systems, the Qwen2.5 series—developed by Alibaba Cloud—establishes a new benchmark for dense parameter scaling. This paper delivers a comprehensive technical and mathematical investigation of Qwen2.5-72B, an open-weights dense language model comprising 72.7 billion parameters across 80 Transformer layers. We analyze its underlying architectural mechanisms, including: (1) Grouped-Query Attention (GQA) operating at an 8:1 query-to-KV head ratio combined with dual-layer QK-Norm, (2) SwiGLU-activated feed-forward networks with expanded intermediate projections, (3) Rotary Position Embedding (RoPE) parameterized with a 1,000,000 base frequency and YaRN context extension enabling 128K token context understanding, and (4) an expanded 152,064-token Byte-Pair Encoding (BPE) vocabulary optimized for multilingual fluency and synthetic reasoning tokens. Furthermore, we investigate its training dynamics across 18 trillion diverse tokens and evaluate its inference serving economics across FP16 and AWQ/GPTQ 4-bit quantization regimes. Our empirical synthesis highlights the architectural refinements that allow Qwen2.5-72B to outperform comparable 70B-class dense models and rival significantly larger proprietary systems across mathematics, coding, and instruction-following benchmarks.
\end{abstract}

\begin{IEEEkeywords}
Qwen2.5-72B, Dense Transformer, Grouped-Query Attention (GQA), QK-Norm, SwiGLU, Long-Context YaRN, LLM Inference.
\end{IEEEkeywords}

\section{Introduction}
The scaling trajectory of foundation models has demonstrated that architectural refinement is as critical as raw parameter volume. While sparsely gated architectures reduce operational training FLOPs, dense architectures remain uniquely reliable, free from load-balancing communication overheads and routing instability. Within the dense domain, the release of Qwen2.5-72B by Alibaba Cloud marks a pivotal milestone in open-weights artificial intelligence \cite{qwen2024qwen25, qwen2024qwen2}.

Pre-trained on a massive 18-trillion token multilingual corpus, Qwen2.5-72B was engineered to overcome historical deficiencies in dense architectures: specifically, quadratic memory scaling during long-sequence decoding, vulnerability to attention score explosion across deep networks, and tokenization inefficiencies across non-Latin scripts and source code.

This paper presents an exhaustive architectural analysis of Qwen2.5-72B. We dissect its structural configuration, derive its mathematical formulations, analyze its long-context extrapolation mechanisms, and examine the memory bandwidth and computational characteristics governing its inference efficiency.

\section{Core Architectural Specifications and Topology}
Qwen2.5-72B is structured as an autoregressive, causal decoder-only Transformer. Its macro-architectural layout comprises 80 identical Transformer blocks operating over a hidden representation dimension $d_{model} = 8,192$. Table \ref{tab:qwen_specs} details the structural hyper-parameters of the model.

\begin{table}[htbp]
\caption{Hyper-parameter Specifications of Qwen2.5-72B}
\label{tab:qwen_specs}
\centering
\begin{tabular}{ll}
\toprule
\textbf{Architectural Hyper-parameter} & \textbf{Value / Specification} \\
\midrule
Total Parameters & 72.7 Billion \\
Number of Layers ($L$) & 80 \\
Hidden Dimension ($d_{model}$) & 8,192 \\
FFN Intermediate Dimension ($d_{ffn}$) & 29,568 \\
Query Attention Heads ($h_Q$) & 64 \\
Key/Value Attention Heads ($h_{KV}$) & 8 \\
Head Dimension ($d_k = d_{model} / h_Q$) & 128 \\
Compression Ratio ($h_Q / h_{KV}$) & 8:1 \\
Vocabulary Size ($|V|$) & 152,064 \\
Native Context Window & 32,768 tokens \\
Extended Context Window & 131,072 tokens (128K) \\
Positional Embedding & RoPE ($\theta_{base} = 1,000,000$) \\
Normalization Scheme & RMSNorm + QK-Norm \\
Activation Function & SwiGLU \\
Pre-training Token Budget & 18.0 Trillion Tokens \\
\bottomrule
\end{tabular}
\end{table}

\section{Attention Architecture: GQA, QK-Norm, and YaRN}

\subsection{Grouped-Query Attention (GQA)}
In standard Multi-Head Attention (MHA) \cite{vaswani2017attention}, every query head possesses a distinct key and value head ($h_Q = h_{KV} = 64$). During autoregressive decoding, loading these unshared KV caches creates severe memory bandwidth saturation. Qwen2.5-72B implements Grouped-Query Attention (GQA) \cite{ainslie2023gqa}, partitioning the 64 query heads into 8 distinct groups, where each group of 8 query heads shares a single Key-Value head pair:
\begin{equation}
g = \frac{h_Q}{h_{KV}} = \frac{64}{8} = 8
\end{equation}
Let $X \in \mathbb{R}^{T \times d_{model}}$ denote the input tensor. The projection operations are defined as:
\begin{align}
Q &= X W_Q, \quad W_Q \in \mathbb{R}^{d_{model} \times (h_Q \cdot d_k)} \\
K &= X W_K, \quad W_K \in \mathbb{R}^{d_{model} \times (h_{KV} \cdot d_k)} \\
V &= X W_V, \quad W_V \in \mathbb{R}^{d_{model} \times (h_{KV} \cdot d_k)}
\end{align}
For the $i$-th query head ($i \in \{1, \dots, 64\}$), its associated Key and Value representations are drawn from group index $j = \lfloor (i-1)/g \rfloor + 1$. This structural sharing reduces the KV cache size by a factor of 8 without degrading linguistic representation fidelity.

\subsection{Mathematical Formulation of QK-Norm}
As Transformers scale in depth to 80 layers, dot-product attention scores $Q K^T / \sqrt{d_k}$ can grow uncontrollably in magnitude, pushing the softmax function into regions of near-zero gradients. This leads to attention entropy collapse and destabilizes mixed-precision (BF16/FP8) training.

To guarantee stability, Qwen2.5-72B incorporates \textbf{QK-Norm}. Root Mean Square Normalization (RMSNorm) \cite{zhang2019rmsnorm} is applied individually to each query and key head prior to computing attention logits:
\begin{equation}
\operatorname{RMSNorm}(z) = \frac{z}{\sqrt{\frac{1}{d_k} \sum_{m=1}^{d_k} z_m^2 + \epsilon}} \odot \gamma
\end{equation}
where $\gamma \in \mathbb{R}^{d_k}$ is a learnable affine parameter. The attention computation for head $i$ becomes:
\begin{equation}
\tilde{Q}_i = \operatorname{RMSNorm}(Q_i), \quad \tilde{K}_j = \operatorname{RMSNorm}(K_j)
\end{equation}
\begin{equation}
\operatorname{Attn}_i = \operatorname{Softmax}\left( \frac{\tilde{Q}_i \tilde{K}_j^T}{\sqrt{d_k}} + M \right) V_j
\end{equation}
where $M$ is the causal masking matrix. Because $\|\tilde{Q}_i\|_2 \approx \sqrt{d_k}$ and $\|\tilde{K}_j\|_2 \approx \sqrt{d_k}$, the dot product is strictly bounded by the Cauchy-Schwarz inequality:
\begin{equation}
\left| \frac{\tilde{Q}_i \tilde{K}_j^T}{\sqrt{d_k}} \right| \le \frac{\|\tilde{Q}_i\|_2 \|\tilde{K}_j\|_2}{\sqrt{d_k}} \approx \sqrt{d_k}
\end{equation}
With $d_k = 128$, the maximum unmasked logit magnitude is bounded by $\approx 11.31$, completely preventing softmax saturation.

\subsection{Rotary Position Embedding (RoPE) and YaRN Scaling}
Qwen2.5-72B encodes positional geometry via Rotary Position Embedding (RoPE) \cite{su2024roformer}. For token position $m$ and vector $x \in \mathbb{R}^{d_k}$, RoPE applies a block-diagonal rotation matrix $R_{\Theta, m}^d$:
\begin{equation}
R_{\Theta, m}^d = \operatorname{diag}\left( R_{\theta_1, m}, R_{\theta_2, m}, \dots, R_{\theta_{d_k/2}, m} \right)
\end{equation}
where $\theta_r = \theta_{base}^{-2(r-1)/d_k}$. Crucially, Qwen2.5 sets $\theta_{base} = 1,000,000$ (compared to standard 10,000 in early Transformers), dampening the rotation frequency of higher dimensions and preventing high-frequency phase collisions across long contexts.

To extrapolate the context window from 32,768 to 131,072 tokens (128K), Qwen2.5 integrates YaRN (Yet another RoPE extensioN) \cite{peng2023yarn}. YaRN applies a piecewise non-uniform frequency scaling factor $s_r$:
\begin{equation}
\theta'_r = (1 - \alpha_r) \frac{\theta_r}{s} + \alpha_r \theta_r
\end{equation}
where $\alpha_r$ smoothly interpolates between linear interpolation at low frequencies (preserving global distance relationships) and zero interpolation at high frequencies (preserving local positional resolution), modulated by an attention temperature factor $t = 0.1 \ln s + 1$.

\section{Feed-Forward Network (FFN) and Tokenization}

\subsection{SwiGLU Activation Function}
The feed-forward block utilizes the Swish-Gated Linear Unit (SwiGLU) \cite{shazeer2020glu}. For an intermediate dimension $d_{ffn} = 29,568$:
\begin{align}
\operatorname{SwiGLU}(x) &= \left( \operatorname{Swish}(x W_{gate}) \odot x W_{up} \right) W_{down} \\
\operatorname{Swish}(z) &= z \cdot \sigma(\beta z) = \frac{z}{1 + e^{-\beta z}}
\end{align}
where $W_{gate}, W_{up} \in \mathbb{R}^{d_{model} \times d_{ffn}}$ and $W_{down} \in \mathbb{R}^{d_{ffn} \times d_{model}}$.
Setting $d_{ffn} = 29,568 \approx \frac{8}{3} d_{model}$ maintains an optimal parameter balance between attention projections and non-linear feature transformation.

\subsection{High-Capacity Multilingual Tokenizer}
Qwen2.5-72B features an expansive vocabulary of $|V| = 152,064$ tokens utilizing a byte-level Byte-Pair Encoding (BPE) tokenizer constructed via tiktoken. Unlike traditional tokenizers biased toward English morphology, Qwen2.5's vocabulary allocates balanced subword capacity across Mandarin, Arabic, Cyrillic, and European languages, as well as programming syntax (Python, C++, Java, JSON, SQL). This high compression ratio reduces sequence lengths by 15--35\% for non-Latin and technical corpora, directly translating to proportional inference speedups.

\begin{figure}[htbp]
\centering
\fbox{\parbox{0.92\linewidth}{
\small
\textbf{Qwen2.5-72B Layer Execution Flow:}\\[4pt]
\textbf{Input Tensor} $X_l \in \mathbb{R}^{T \times 8192}$ \\
$\quad \downarrow$ \\
\textbf{Pre-RMSNorm} $\to X_{norm} = \operatorname{RMSNorm}(X_l)$ \\
$\quad \downarrow$ \\
\textbf{GQA Module (64 Q Heads, 8 KV Heads):} \\
$\quad \bullet$ Projections: $Q = X_{norm} W_Q$, $K = X_{norm} W_K$, $V = X_{norm} W_V$ \\
$\quad \bullet$ QK-Norm: $\tilde{Q} = \operatorname{RMSNorm}(Q)$, $\tilde{K} = \operatorname{RMSNorm}(K)$ \\
$\quad \bullet$ Positional Rotation: $Q_{rot} = \operatorname{RoPE}(\tilde{Q}, \theta=10^6)$, $K_{rot} = \operatorname{RoPE}(\tilde{K}, \theta=10^6)$ \\
$\quad \bullet$ Attn: $\operatorname{Softmax}\left(\frac{Q_{rot} K_{rot}^T}{\sqrt{128}} + M\right) V$ \\
$\quad \bullet$ Output Projection: $Y_{attn} = \operatorname{Attn} \cdot W_O$ \\
$\quad \downarrow$ (Residual Addition) \\
$X_{mid} = X_l + Y_{attn}$ \\
$\quad \downarrow$ \\
\textbf{Pre-RMSNorm} $\to X'_{norm} = \operatorname{RMSNorm}(X_{mid})$ \\
$\quad \downarrow$ \\
\textbf{SwiGLU FFN Block ($d_{ffn} = 29,568$):} \\
$\quad \bullet Y_{ffn} = (\operatorname{Swish}(X'_{norm} W_{gate}) \odot X'_{norm} W_{up}) W_{down}$ \\
$\quad \downarrow$ (Residual Addition) \\
\textbf{Output Tensor} $X_{l+1} = X_{mid} + Y_{ffn}$
}}
\caption{Architectural structure of a single Qwen2.5-72B Transformer layer.}
\label{fig:qwen_layer}
\end{figure}

\section{Training Strategy and Optimization Dynamics}

\subsection{Pre-training Corpus and Compute Budget}
Qwen2.5-72B was trained on 18 trillion tokens—a massive scale among open-weights dense models. The training mixture prioritized synthetic math generation, code documentation, and multi-step reasoning traces. Data deduplication, perplexity filtering, and synthetic instruction inversion were applied to maximize data density.

\subsection{Optimization Setup}
Pre-training utilized the AdamW optimizer with $\beta_1 = 0.9$, $\beta_2 = 0.95$, and a weight decay coefficient of $0.1$. The learning rate followed a cosine annealing schedule with a linear warmup over the first 2,000 steps, peaking at $\eta_{max} = 1.2 \times 10^{-4}$ before decaying to $1.2 \times 10^{-5}$. Numerical stability across 80 layers was preserved via mixed-precision BF16 computations coupled with gradient clipping at a norm of $1.0$.

\subsection{Post-Training: Alignment and Preference Optimization}
The post-training regime utilized a multi-stage methodology:
\begin{enumerate}
    \item \textbf{Supervised Fine-Tuning (SFT):} Conducted over 1 million diverse instructional demonstrations focusing on structured JSON formatting, function calling, mathematical derivation, and agentic workflows.
    \item \textbf{Direct Preference Optimization (DPO):} Aligned model responses with human preferences without learning a separate reward network \cite{rafailov2024direct}.
    \item \textbf{Online RLHF / GRPO:} Applied rule-based verifiers for mathematical correctness and unit-test execution for code synthesis, maximizing objective reasoning fidelity.
\end{enumerate}

\section{Inference Efficiency and Serving Economics}

\subsection{KV Cache Memory Footprint}
For an autoregressive model, the KV cache per token per layer is given by:
\begin{equation}
\text{KV}_{\text{token}} = 2 \times h_{KV} \times d_k \times \text{bytes}
\end{equation}
Under FP16 (2 bytes per parameter), with $h_{KV} = 8$ and $d_k = 128$:
\begin{equation}
\text{KV}_{\text{token/layer}} = 2 \times 8 \times 128 \times 2 = 4,096 \text{ bytes} = 4 \text{ KB}
\end{equation}
Across all $L = 80$ layers, the cache requirement is:
\begin{equation}
\text{Total KV}_{\text{token}} = 80 \times 4 \text{ KB} = 320 \text{ KB/token}
\end{equation}
In contrast, if standard Multi-Head Attention ($h_{KV} = 64$) were used, the requirement would be $8 \times 320 \text{ KB} = 2.56 \text{ MB/token}$. At the maximum context length of 128,000 tokens:
\begin{itemize}
    \item \textbf{MHA 128K Cache:} $128,000 \times 2.56 \text{ MB} \approx 327.68 \text{ GB}$ (Exceeds 4x A100 80GB VRAM solely for cache).
    \item \textbf{Qwen2.5 GQA 128K Cache:} $128,000 \times 320 \text{ KB} \approx 40.96 \text{ GB}$ (Easily fits within distributed multi-GPU serving environments).
\end{itemize}

\subsection{Quantization Regimes: AWQ and GPTQ}
Serving 72.7B parameters in native BF16 requires:
\begin{equation}
\text{VRAM}_{\text{weights}} = 72.7 \times 10^9 \times 2 \approx 145.4 \text{ GB}
\end{equation}
necessitating at least two 80GB GPUs (Tensor Parallelism $TP=2$). Through Activation-aware Weight Quantization (AWQ) \cite{lin2023awq} or GPTQ 4-bit quantization, the weight footprint drops to:
\begin{equation}
\text{VRAM}_{\text{INT4}} \approx \frac{145.4 \text{ GB}}{4} + \text{overheads} \approx 42.5 \text{ GB}
\end{equation}
This allows full-parameter inference of a 72B-class frontier model on a single 48GB GPU (e.g., NVIDIA RTX 6000 Ada or A6000) or two consumer 24GB GPUs (RTX 4090).

\section{Benchmark Performance and Comparative Analysis}
Table \ref{tab:benchmarks} compares Qwen2.5-72B-Instruct against prominent contemporary open-weights models, including Llama 3.1 70B \cite{dubey2024llama3} and DeepSeek-V2.5.

\begin{table}[htbp]
\caption{Empirical Benchmark Evaluation Across Frontier Models}
\label{tab:benchmarks}
\centering
\begin{tabular}{lccc}
\toprule
\textbf{Benchmark / Metric} & \textbf{Qwen2.5-72B} & \textbf{Llama 3.1 70B} & \textbf{DeepSeek-V2.5} \\
\midrule
MMLU (5-shot) & \textbf{86.1\%} & 83.6\% & 80.6\% \\
MMLU-Pro (0-shot CoT) & \textbf{66.4\%} & 56.2\% & 51.6\% \\
GSM8K (8-shot CoT) & \textbf{91.6\%} & 88.9\% & 89.2\% \\
MATH (0-shot CoT) & \textbf{59.8\%} & 42.2\% & 43.4\% \\
HumanEval (pass@1) & \textbf{86.4\%} & 80.5\% & 81.1\% \\
LiveCodeBench & \textbf{41.8\%} & 34.4\% & 35.8\% \\
Arena-Hard & \textbf{81.2} & 69.3 & 76.2 \\
Context Window & 128K & 128K & 128K \\
Active Parameters & 72.7B & 70.6B & 21B (MoE) \\
\bottomrule
\end{tabular}
\end{table}

The benchmark results reflect the benefits of Qwen2.5's architectural choices:
\begin{enumerate}
    \item \textbf{Mathematical Superiority:} Scoring 59.8\% on MATH (a +17.6\% gain over Llama 3.1 70B) demonstrates the efficacy of QK-Norm stabilized pre-training combined with heavy synthetic STEM token exposure.
    \item \textbf{Code Generation:} HumanEval pass@1 of 86.4\% confirms the utility of its 152K-token vocabulary in preserving compact code syntax without morphological fragmentation.
\end{enumerate}

\section{Conclusion}
Qwen2.5-72B demonstrates the immense power of disciplined dense Transformer engineering. Through the synergy of Grouped-Query Attention, QK-Norm stabilization, 1M-base RoPE with YaRN context extension, and a high-capacity multilingual tokenizer, Alibaba Cloud has produced an open-weights dense architecture that defines state-of-the-art capability in mathematics, coding, and long-context comprehension.

\begin{thebibliography}{99}

\bibitem{qwen2024qwen25}
Qwen Team, A.~Yang, \textit{et al.}, ``Qwen2.5 technical report,'' \textit{arXiv preprint arXiv:2412.15115}, 2024.

\bibitem{qwen2024qwen2}
A.~Yang, B.~Yang, B.~Hui, B.~Zheng, B.~Yu, C.~Zhou, C.~Li, C.~Li, D.~Liu, F.~Huang, \textit{et al.}, ``Qwen2 technical report,'' \textit{arXiv preprint arXiv:2407.10671}, 2024.

\bibitem{vaswani2017attention}
A.~Vaswani, N.~Shazeer, N.~Parmar, J.~Uszkoreit, L.~Jones, A.~N.~Gomez, {\L}.~Kaiser, and I.~Polosukhin, ``Attention is all you need,'' in \textit{Advances in Neural Information Processing Systems (NeurIPS)}, vol.~30, 2017.

\bibitem{ainslie2023gqa}
J.~Ainslie, J.~Lee-Thorp, M.~de~Jong, Y.~Zemlyanskiy, F.~Lebr{\'o}n, and S.~Sanghai, ``GQA: Training generalized multi-query transformer models from multi-head checkpoints,'' \textit{arXiv preprint arXiv:2305.13245}, 2023.

\bibitem{zhang2019rmsnorm}
B.~Zhang and R.~Sennrich, ``Root mean square layer normalization,'' in \textit{Advances in Neural Information Processing Systems (NeurIPS)}, vol.~32, 2019.

\bibitem{su2024roformer}
J.~Su, M.~Ahmed, Y.~Lu, S.~Pan, W.~Bo, and Y.~Liu, ``RoFormer: Enhanced transformer with rotary position embedding,'' \textit{Neurocomputing}, vol.~568, p.~127063, 2024.

\bibitem{peng2023yarn}
B.~Peng, J.~Quesnelle, H.~Fan, and E.~Shippole, ``YaRN: Efficient context window extension of large language models,'' \textit{arXiv preprint arXiv:2309.00071}, 2023.

\bibitem{shazeer2020glu}
N.~Shazeer, ``GLU variants improve transformer,'' \textit{arXiv preprint arXiv:2002.05202}, 2020.

\bibitem{lin2023awq}
J.~Lin, J.~Tang, H.~Tang, S.~Yang, X.~Dang, and S.~Han, ``AWQ: Activation-aware weight quantization for LLM compression and acceleration,'' \textit{arXiv preprint arXiv:2306.00978}, 2023.

\bibitem{rafailov2024direct}
R.~Rafailov, A.~Sharma, E.~Mitchell, S.~Ermon, C.~D.~Manning, and C.~Finn, ``Direct preference optimization: Your language model is secretly a reward model,'' in \textit{Advances in Neural Information Processing Systems (NeurIPS)}, vol.~36, 2023.

\bibitem{dubey2024llama3}
A.~Dubey, A.~Jauhri, A.~Pandey, A.~Kadian, A.~Al-Dahle, A.~Letman, A.~Mathur, A.~Schelten, A.~Yang, A.~Fan, \textit{et al.}, ``The Llama 3 herd of models,'' \textit{arXiv preprint arXiv:2407.21783}, 2024.

\end{thebibliography}

\end{document}
